% ------------------------------------------------------------------------------
% 1 Introduction
% ------------------------------------------------------------------------------
\section{Introduction}\label{sec:intro}

Colorectal cancer is one of the most common malignancies worldwide, and
for patients undergoing potentially curative resection the question of
whether systemic adjuvant chemotherapy should follow surgery remained
open for much of the twentieth century. The North Central Cancer
Treatment Group (NCCTG) program addressed this question directly:
building on an earlier randomized comparison of observation against
levamisole alone \citep{laurie1989}, the investigators added a third arm
combining the immunomodulator levamisole with the antimetabolite
5-fluorouracil (5-FU) \citep{moertel1990}. The combination reduced
recurrence and mortality substantially and became a standard of care for
resected stage B/C colon carcinoma. The individual-patient data from
this program were subsequently contributed to the \texttt{survival}
package for R \citep{therneau2026}, where they now serve as one of the
canonical datasets for teaching survival analysis; the version analysed
here contains 929 patients from a single study.

Re-analyzing an old trial with new methods is more than an academic
exercise, and five considerations motivated each component of this
project. First, endpoint definitions and analysis conventions have
evolved: modern adjuvant colorectal trials define disease-free survival
(DFS) as the time to recurrence \emph{or} death, whereas the published
coding of this dataset treats death without recurrence as censoring; a
rigorous re-analysis must therefore address both views and, importantly,
the competing-risks consequences of the difference. Second, missing-data
handling, covariate-adjusted estimation, and absolute-benefit reporting
have all been refined since the original publications, and regulatory
guidance now actively encourages pre-specified covariate adjustment in
the primary analysis of randomized trials
\citep{fda2023,kahan2014,sterne2009}. Third, the trial-design landscape
has changed: group-sequential designs with alpha- and beta-spending
functions are now the default in oncology, and it is natural to ask
whether this trial---run with a fixed-sample mindset---would have
stopped early under a modern monitoring plan
\citep{landemets1983,o1979}. Fourth, machine-learning prediction models
and Bayesian machinery offer perspectives on prognosis and evidence
accumulation that complement the classical toolkit
\citep{harrell2015,gelman2013}. Finally, the modern estimand framework
and recent methods for heterogeneous treatment effects invite a
re-examination of \emph{for whom} the benefit is largest
\citep{hernan2020}.

\subsection*{Contributions}

This paper integrates all of these perspectives into a single coherent,
fully scripted pipeline and makes the following contributions:

\begin{enumerate}[leftmargin=1.7em, itemsep=2pt, topsep=3pt]
  \item \textbf{A comprehensive treatment-effect re-estimate} on the
        composite DFS endpoint, combining multiple imputation,
        proportional-hazards diagnostics, competing-risk (Aalen--Johansen
        and Fine--Gray) sensitivity analyses, and absolute-benefit
        measures---five-year risk differences, numbers needed to treat,
        and restricted mean survival times (RMST), the latter providing a
        collapsible, assumption-light summary that directly contrasts
        with the non-collapsibility of the hazard ratio
        \citep{royston2011,uno2014}.
  \item \textbf{A counterfactual trial redesign and monitoring
        re-enactment}: the trial is re-sized and re-designed as a
        three-look group-sequential study
        \citep{o1979,landemets1983,hwang1990}, and the monitoring plan is
        replayed on the real event stream, showing that a modern design
        would have stopped the trial for efficacy at the first interim
        analysis, with conditional-power and Bayesian monitoring
        cross-checks.
  \item \textbf{Simulation-calibrated operating characteristics} for the
        design and analysis choices: the efficiency gain from covariate
        adjustment (including a direct quantitative demonstration of
        hazard-ratio non-collapsibility), the type-I error and stopping
        probabilities of the sequential design, and confidence-interval
        coverage under non-proportional hazards.
  \item \textbf{An honest validation study of five survival prediction
        models} spanning penalized regression, kernel ensembles, and
        gradient boosting, quantifying the optimism gap between apparent
        and cross-validated discrimination.
  \item \textbf{A Bayesian re-analysis} with explicit prior sensitivity
        (skeptical, weakly-informative, enthusiastic), event-driven
        sequential monitoring, posterior predictive checks, and
        baseline-model criticism by approximate leave-one-out
        cross-validation \citep{vehtari2017}.
  \item \textbf{A heterogeneous-treatment-effect analysis} on the
        absolute five-year risk scale, combining jackknife
        pseudo-observation regression \citep{andersen2003,parner2010},
        cross-fitted G-computation, and doubly robust AIPW estimation
        \citep{hernan2020}, which locates the largest absolute benefits
        in the highest-risk patients.
\end{enumerate}

Beyond the specific findings, the paper is deliberately
methods-integrative: every paradigm is applied to the same data, under a
common adjustment set and common software-engineering discipline (fixed
seeds, file-based module interfaces, one-command reproduction), so that
the contrasts between paradigms---collapsible versus non-collapsible
estimands, alpha-spending versus posterior monitoring, apparent versus
honest prediction accuracy---are attributable to the methods rather than
to implementation differences.

The remainder of the paper is organized as follows.
Section~\ref{sec:data} describes the data and endpoint construction;
Section~\ref{sec:methods} details all statistical methods;
Section~\ref{sec:results} presents results in five thematic blocks; and
Section~\ref{sec:discussion} discusses strengths, limitations, and
conclusions.
